\section{Residential and commercial stationary combustion}
\label{sec:rescom}

\subsection{Scope and inventory mapping}

Residential and commercial combustion covers small stationary sources that are
not public power or industrial plants: space and water heating, cooking, and
related uses in dwellings and in commercial or institutional buildings. In the
CAMS-REG-ANT GNFR split, these emissions are reported under
\textbf{GNFR~C} (\emph{Other stationary combustion}), together with other
non-industrial combustion that is not mapped to GNFR~A or~B
\citep{cams_glob_ant}. The downscaling workflow described here targets the
\textbf{area-source} component of GNFR~C (CAMS \texttt{source\_type} area),
because that is the part of the inventory that requires redistribution from
coarse CAMS grid cells to the national fine reference grid. Point sources, if
present within the same GNFR code, are out of scope for this block.

The method follows the EMEP/EEA air pollutant emission inventory guidebook
logic for residential and commercial combustion \citep{EMEP_EEA}: emissions
depend on fuel, appliance technology, and activity, while fine-scale spatial
patterns are driven by building-related heat demand and settlement structure
rather than by a single land-cover class.

\subsection{Conceptual design}

GNFR~C aggregates many fuels and devices into one sectoral total per CAMS cell
and pollutant. A single land-cover proxy would force all pollutants to share
the same spatial pattern. The upgraded approach instead combines:
\begin{itemize}
  \item \textbf{National activity structure} from GAINS domestic-sector shares
        (fuel $\times$ appliance, year 2020), mapped to a small set of model
        classes;
  \item \textbf{Technology-specific emission factors} from EMEP/EEA-style
        tabulations (fuel and appliance categories), yielding class- and
        pollutant-dependent emission intensities;
  \item \textbf{End-use calibration} from Eurostat household final energy by
        end use, used as a country-specific multiplicative correction on mapped
        activity;
  \item \textbf{Spatial proxies} from Hotmaps building heat and floor-area
        densities, refined with CORINE Land Cover (CLC) urban morphology
        classes.
\end{itemize}
Pollutant-specific downscaling weights are built inside each CAMS cell so that
fine-grid allocations sum to the CAMS cell total per pollutant (mass
conserving).

\subsection{Model classes}

Process detail is represented through seven aggregated \emph{model classes}
$k$: five residential ($R_{\ast}$) and two commercial ($C_{\ast}$), listed in
Appendix~\ref{tab:rescom_model_classes}. They condense GAINS ``fuel $\times$
device'' rows: each GAINS row is assigned to exactly one class using ordered
rules on the GAINS fuel string and appliance string (first match wins). Rules
are fuel-aware where needed so that the same appliance phrase in different
fuels can still be aligned with the correct EMEP fuel category via explicit
hints (Appendix~\ref{tab:rescom_emep_fuel_hints}).

\subsection{National activity and end-use correction}

Let $r$ denote the country (ISO3) associated with a CAMS source record, $(f,a)$
a GAINS fuel--appliance pair, and $s_{r,f,a}$ the national share of fuel $f$
allocated to appliance $a$ in the chosen GAINS year (converted to a fraction).
Let $k(f,a)$ be the model class returned by the mapping rules. The
end-use correction $F_{\mathrm{enduse}}(r,k)$ is derived from Eurostat household
final energy shares by end use (Table~3 in the custom workbook; exact row
labels in Appendix~\ref{tab:rescom_eurostat_table3}). Residential classes
linked to space and water heating use the mean of the corresponding Eurostat
shares where both are available. Commercial boiler classes use a neutral
factor (unity) because the same household end-use table does not describe
commercial buildings. If Eurostat parsing fails for a country, all factors
default to unity.

The corrected activity weight for row $(f,a)$ is
\begin{equation}
  A_{\mathrm{corr}}(r,f,a) = s_{r,f,a}\, F_{\mathrm{enduse}}\bigl(r,k(f,a)\bigr).
\end{equation}

\subsection{Emission factors and class intensities}

For each pollutant $p$, emission factors $\mathrm{EF}_p(f,a)$ are taken from
technology tables keyed by fuel and appliance character strings consistent with
the EMEP/EEA guidebook (chapter~1A4 residential and commercial combustion).
The implementation matches each GAINS row to the best-scoring table row using
token overlap on fuel and appliance text, with an optional bonus when a
user-defined hint links a GAINS fuel substring to an EMEP fuel substring
(Appendix~\ref{tab:rescom_emep_fuel_hints}).

The class-level intensity used in the weight builder is
\begin{equation}
  M(r,p,k) = \sum_{(f,a):\, k(f,a)=k} A_{\mathrm{corr}}(r,f,a)\,
             \mathrm{EF}_p(f,a).
\end{equation}
Thus $M$ combines national structure, end-use calibration, and guidebook
factors before any spatial pattern is applied. Pollutant-specific weights arise
because $\mathrm{EF}_p$ differs across $p$ (including separate treatment of
fossil and biogenic CO$_2$ when the inventory provides \texttt{co2\_ff} and
\texttt{co2\_bf}; total CO$_2$ uses their sum when configured).

\subsection{Spatial base proxies and CLC morphology}

Let $i$ index pixels on the national CORINE reference window (same extent and
resolution as the SourceProxies outputs: CLC grid aligned to the NUTS-based
window). Hotmaps rasters supply residential and non-residential current heat
density and gross floor area density, denoted
$H_{\mathrm{res}}(i)$, $H_{\mathrm{com}}(i)$,
$G_{\mathrm{res}}(i)$, $G_{\mathrm{com}}(i)$, harmonised to that grid.

Base proxies combine heat and floor area with exponents
$\gamma_{\mathrm{heat}}=0.7$ and $\gamma_{\mathrm{gfa}}=0.3$, or an additive
alternative if a multiplicative form is numerically unstable:
\begin{align}
  R_{\mathrm{base}}(i) &=
  H_{\mathrm{res}}(i)^{\gamma_{\mathrm{heat}}}\,
  G_{\mathrm{res}}(i)^{\gamma_{\mathrm{gfa}}},
  \\
  C_{\mathrm{base}}(i) &=
  H_{\mathrm{com}}(i)^{\gamma_{\mathrm{heat}}}\,
  G_{\mathrm{com}}(i)^{\gamma_{\mathrm{gfa}}},
\end{align}
with non-negative inputs and missing values set to zero after resampling.

CLC classes 111 (continuous urban fabric), 112 (discontinuous urban fabric), and
121 (industrial or commercial units) define binary indicators
$U_{111}(i)$, $U_{112}(i)$, $U_{121}(i)$.
For residential fireplace and heating-stove classes, a morphology factor
$f_R(i)$ upweights discontinuous urban fabric and downweights purely industrial
cells relative to other locations; for commercial boiler classes, a factor
$f_C(i)$ emphasises industrial/commercial and dense urban cells. Cooking stoves
and residential boilers without that morphology use $R_{\mathrm{base}}(i)$ only.
Full algebraic forms are given in the configuration; the role of CLC is to
modulate \emph{where} within the building-energy proxy residential versus
commercial combustion is emphasised, not to replace energy-based allocation.

The spatial score for class $k$ is denoted $X(i,k)$: piecewise products of
$R_{\mathrm{base}}$ or $C_{\mathrm{base}}$ with $f_R$ or $f_C$ as defined for
each $k$.

\subsection{Unnormalised weights, normalisation, and gridded emissions}

For pollutant $p$ and fine pixel $i$ belonging to CAMS area cell $c$, let
$r(i)$ be the country of the inventory record (taken from the CAMS country
index for that cell). The unnormalised weight is
\begin{equation}
  U(i,p) = \sum_{k} M\bigl(r(i),p,k\bigr)\, X(i,k).
\end{equation}
Within each CAMS cell $c$,
\begin{equation}
  w(i,p) = \frac{U(i,p)}{\displaystyle\sum_{j \in c} U(j,p)},
\end{equation}
with the convention that if the denominator is zero or non-finite, weights are
uniform over pixels intersecting $c$ inside the national domain. By construction,
$\sum_{i \in c} w(i,p)=1$ for each $p$.

Gridded emissions (kg\,yr$^{-1}$ per pixel) follow
\begin{equation}
  E(i,p) = w(i,p)\, E_{\mathrm{CAMS}}(c,p),
\end{equation}
where $E_{\mathrm{CAMS}}(c,p)$ is the CAMS GNFR~C area total for cell $c$ and
pollutant $p$. The same workflow applies to NH$_3$, NO$_x$, SO$_x$, NMVOC,
PM$_{2.5}$, PM$_{10}$, CO, and CO$_2$ (as total fossil plus biogenic CO$_2$
when both components are summed).

\subsection{Inputs and software}

Table~\ref{tab:rescom_inputs} summarises mandatory inputs; Eurostat Table~1 fuel
columns (Appendix~\ref{tab:rescom_eurostat_table1}) are documented for
consistency with national fuel mix but are not yet used to rescale $M$.
Implementation lives under \path{Residential/downscale/} with JSON
configuration (\path{Residential/config/residential_cams_downscale.config.json},
\path{Residential/config/eurostat_end_use.json},
\path{Residential/config/GAINS_mapping.json},
\path{Residential/config/EMEP_emission_factors.json}). Hotmaps GeoTIFFs can be
retrieved with \path{Residential/auxiliaries/download_hotmaps_building_rasters.py}.

\begin{table}[htbp]
  \centering
  \caption{Main inputs for GNFR~C residential/commercial downscaling.}
  \label{tab:rescom_inputs}
  \begin{tabular}{p{3.2cm}p{9.5cm}}
    \hline
    Ingredient & Role \\
    \hline
    CAMS-REG-ANT NetCDF & GNFR~C area emissions $E_{\mathrm{CAMS}}(c,p)$,
      cell geometry, country index \\
    CORINE CLC 2018 & Reference grid, classes 111/112/121 for morphology \\
    Hotmaps (heat \& GFA) & $H_{\ast}$, $G_{\ast}$ building proxies \\
    GAINS \texttt{dom\_share\_ENE\_*} & National fuel--appliance shares \\
    Eurostat custom workbook & $F_{\mathrm{enduse}}$ from household end-use
      shares (Table~3) \\
    EMEP-style EF tables & $\mathrm{EF}_p(f,a)$ by fuel and appliance \\
    NUTS geometry & National bounding window for the CLC reference grid \\
    \hline
  \end{tabular}
\end{table}

\subsection{Limitations}

\paragraph{Sector purity.} GNFR~C may include other stationary combustion than
dwellings and services; the proxy cannot separate processes that CAMS does not
split.

\paragraph{Uniform national GAINS structure.} Appliance shares are applied
uniformly within each country; there is no subnational GAINS disaggregation.

\paragraph{Eurostat and commercial classes.} Household end-use factors do not
describe commercial energy use; commercial classes rely on unity end-use
factors unless extended later.

\paragraph{EF matching.} Automatic string matching between GAINS and EMEP rows
can mis-assign rare fuels; the hint table and manual EF tables should be
audited for the countries of interest.

\paragraph{CAMS cell versus national window.} Pixels outside the national CLC
window or CAMS cells that do not intersect the window receive no allocation from
those cells; boundary effects follow the same logic as for other CORINE-based
proxies.

Further tabulated detail (model classes, Eurostat labels, EMEP fuel hints)
appears in Appendix~\ref{app:rescom_downscaling}.
